\chapter{Сферы, гиперболическое пространство и смысл \(k\)}

Давайте немного займёмся математикой — простой математикой. Сферы в любом числе измерений — особые геометрии: они равномерны, и здесь равномерность означает постоянную кривизну. Сначала идёт положительная кривизна, затем равномерная отрицательная кривизна, а после этого — связь с уравнениями Эйнштейна.

Сферическая вселенная означает, что сферично \emph{пространство}, а не пространство-время. Есть и другая возможность — равномерная отрицательная кривизна. Именно сопоставление этих случаев делает ясной роль общей теории относительности.

\section{Сферы и равномерная геометрия}

1-сфера — это просто окружность. Проще всего описать её, вложив в двумерное вспомогательное пространство и записав
\begin{equation}
x^2+y^2=R^2.
\end{equation}
Единичная 1-сфера задаётся так:
\begin{equation}
x^2+y^2=1.
\end{equation}

2-сфера описывается точно так же, только теперь используются три вспомогательные координаты:
\begin{equation}
x^2+y^2+z^2=R^2,
\end{equation}
в единичном случае
\begin{equation}
x^2+y^2+z^2=1.
\end{equation}

Эти координаты — лишь вспомогательный приём. Интересующая нас геометрия — не окружающее пространство \(x,y,z\). Речь идёт о самой поверхности. Существо, живущее на 1-сфере, живёт на окружности; у него нет доступа ни к внутренней области, ни к внешней. Та же идея относится и к 2-сфере. Заполнение сферы и пространство вне неё принадлежат только вспомогательной картине.

Та же конструкция даёт 3-сферу. Вводят ещё одну координату, скажем \(w\), и пишут
\begin{equation}
x^2+y^2+z^2+w^2=R^2,
\end{equation}
а в единичном случае
\begin{equation}
x^2+y^2+z^2+w^2=1.
\end{equation}
3-сферу представить труднее, потому что теперь вспомогательное пространство было бы четырёхмерным, но сама математика не намного сложнее.

Стандартное обозначение — \(\Omega_n\) для \(n\)-сферы. Поэтому \(\Omega_1\) — это 1-сфера, \(\Omega_2\) — 2-сфера, а \(\Omega_3\) — 3-сфера.

\section{Метрики на сфере}

Более полезно внутреннее описание. Сферу описывают координатами, живущими на ней самой.

Для единичной 1-сферы достаточно одной координаты \(\phi\), причём \(0\leq \phi \leq 2\pi\). Её метрика имеет вид
\begin{equation}
ds^2=d\phi^2.
\end{equation}
Это также записывают как
\begin{equation}
d\Omega_1^2=d\phi^2.
\end{equation}

Единичную 2-сферу можно описать как семейство 1-сфер. Один угол \(\phi\) по-прежнему обходит каждую окружность, а второй угол \(\theta\) идёт от одного конца сферы до другого. Этот второй угол изменяется от \(0\) до \(\pi\). Движение в направлении \(\theta\) даёт вклад \(d\theta^2\). Движение в направлении \(\phi\) происходит по окружности радиуса \(\sin\theta\). Поэтому
\begin{equation}
d\Omega_2^2=d\theta^2+\sin^2\theta\,d\phi^2
\end{equation}
или
\begin{equation}
d\Omega_2^2=d\theta^2+\sin^2\theta\,d\Omega_1^2.
\end{equation}

3-сфера получается тем же способом, теперь как семейство 2-сфер. Введём ещё одну координату \(\alpha\). Тогда
\begin{equation}
d\Omega_3^2=d\alpha^2+\sin^2\alpha\,d\Omega_2^2.
\end{equation}
В развернутом виде это
\begin{equation}
d\Omega_3^2=d\alpha^2+\sin^2\alpha\left(d\theta^2+\sin^2\theta\,d\phi^2\right).
\end{equation}
Здесь \(\alpha\) пробегает от одного конца, где размер равен нулю, до другого, так что \(0\leq \alpha \leq \pi\).\editorialnote{В 00:17:24 в аудиозаписи на мгновение слышится ``от нуля до двух пи'', но окружающая конструкция и явный множитель \(\sin^2\alpha\) требуют, чтобы новая координата пробегала от одного конца семейства сфер меньшей размерности, где размер равен нулю, до другого.}

Общий рекурсивный вид таков:
\begin{equation}
d\Omega_n^2=d\alpha^2+\sin^2\alpha\,d\Omega_{n-1}^2.
\end{equation}

Как только принять эту точку зрения, о картине вложения можно забыть. 1-сфера описывается одной координатой, 2-сфера — двумя, 3-сфера — тремя и так далее. Геометрия существует сама по себе.

Эта геометрия особая, потому что она однородна. Можно было бы деформировать метрику, заменив \(\sin^2\theta\) какой-нибудь другой функцией, которая по-прежнему обращается в нуль на обоих концах. Тогда получилась бы некоторая деформированная замкнутая поверхность, например эллипсоид. Но она уже не была бы равномерной. Сфера однородна, потому что каждая точка на ней такая же, как любая другая.

Если нужна сфера радиуса \(A\), а не единичного радиуса, то единичную метрику умножают на \(A^2\):
\begin{equation}
ds^2=A^2\,d\Omega_3^2.
\end{equation}
Квадрат необходим, потому что метрика измеряет квадрат расстояния.

Короткое отступление касается 0-сферы. Экстраполируя на один шаг назад, получаем
\begin{equation}
x^2=1,
\end{equation}
у которого решениями являются две точки \(x=\pm 1\). Так что единичная 0-сфера — это просто две точки.

\section{Сферическая FRW-вселенная}

Теперь нас интересует не одно только пространство, а пространство-время. Соответствующая картина — это 3-сфера, радиус которой меняется со временем.

На всём протяжении важно одно предупреждение: время отлично от пространственных координат. Не следует представлять расширяющуюся Вселенную как сферу более высокой размерности, в которой время играет роль ещё одной координаты вложения.

Имея это в виду, интервал пространства-времени равен
\begin{equation}
d\tau^2=dt^2-a(t)^2\,d\Omega_3^2.
\end{equation}
Это метрика FRW для вселенной со сферической пространственной геометрией. Единственное новое обстоятельство по сравнению со статической сферой состоит в том, что радиус стал зависящей от времени функцией \(a(t)\).

Удобно выбрать координаты так, чтобы мы находились на одном из полюсов. В этой точке нет ничего особенного; это просто соглашение. Тогда \(\alpha\) можно рассматривать как радиальную координату, отсчитываемую от нас на единичной сфере. Это не радиальное расстояние в километрах. Это безразмерная угловая мера на единичной сфере, но в геометрии она играет роль радиального расстояния.

Если не считать локальных сгустков, таких как галактики и звёзды, то это и есть вся геометрия FRW-вселенной со сферической пространственной геометрией. Её пространственная кривизна положительна, и соответствующее соглашение о знаке записывается так:
\begin{equation}
k=+1.
\end{equation}
Само по себе здесь ничто не равно ``плюс одному''. Это просто означает положительную пространственную кривизну.

\section{Взгляд со сферы}

Как выглядит такая сфера, если смотреть с одной точки на ней? Здесь и появляется стереографическая проекция.

Возьмём 2-сферу, касающуюся бесконечной плоскости в южном полюсе. Выберем точку на сфере. Проведём прямую от северного полюса через эту точку и продолжим её, пока она не пересечёт плоскость. Так каждой точке на сфере ставится в соответствие точка на плоскости. Южный полюс переходит в начало координат на плоскости. Северный полюс — исключительная точка: при приближении к нему с разных направлений его образ уходит на бесконечность, причём по-разному в зависимости от направления. Поэтому его рассматривают как точку на бесконечности.

Если одинаковые ``континенты'' размещать в разных местах на сфере, то их изображения на плоскости становятся больше по мере приближения к северному полюсу. Поэтому наблюдатель, помещённый в южный полюс, видел бы дальние объекты под большими углами, чем в плоском пространстве. При положительной кривизне объект известного размера выглядит слишком большим по сравнению с евклидовой геометрией.

В этом и состоит основная наблюдательная идея. О кривизне судят по связи между расстоянием, размером и угловым размером. Сами галактики не очень хороши в качестве стандартных линеек на очень больших расстояниях, но логика именно такова.

\begin{classroomqa}
\lecturetimestamp{00:45:54}
\audiencequestion{Нужно ли также знать расстояние, а не только угловой размер?}
\lecturerresponse{Да. Одного угла недостаточно. Нужны и угол, и расстояние.}
\end{classroomqa}

\begin{classroomqa}
\lecturetimestamp{00:46:18}
\audiencequestion{Проецируются ли окружности на поверхности сферы в окружности?}
\lecturerresponse{Да. Окружности проецируются в окружности. Это одно из особых свойств стереографической проекции.}
\end{classroomqa}

\begin{classroomqa}
\lecturetimestamp{00:47:02}
\audiencequestion{Сохранил бы континент ту же форму при такой проекции?}
\lecturerresponse{Здесь есть два разных утверждения. Одно состоит в том, что окружности проецируются в окружности. Другое — более слабое — состоит в том, что достаточно малые фигуры сохраняют свою локальную форму, потому что углы сохраняются. Это отображение конформно: оно сохраняет углы, хотя и не относительные размеры.}
\end{classroomqa}

\section{Равномерная отрицательная кривизна}

Случай равномерной отрицательной кривизны вводится через его метрику.

Аналог сферы с отрицательной кривизной — гиперболическое пространство. В том же рекурсивном духе, что и раньше, его метрика получается заменой \(\sin\alpha\) на \(\sinh\alpha\):
\begin{equation}
d\mathcal H_n^2=d\alpha^2+\sinh^2\alpha\,d\Omega_{n-1}^2.
\end{equation}
Угловая часть меньшей размерности остаётся метрикой единичной сферы \(d\Omega_{n-1}^2\).

\begin{classroomqa}
\lecturetimestamp{00:55:51}
\audiencequestion{В метрике этот дифференциальный элемент меньшей размерности по-прежнему сферический?}
\lecturerresponse{Да. Этот фрагмент меньшей размерности — обычная метрика единичной сферы. Меняется множитель \(\sinh^2\alpha\) перед ним.}
\end{classroomqa}

Гиперболический синус равен
\begin{equation}
\sinh\alpha=\frac{e^\alpha-e^{-\alpha}}{2}.
\end{equation}
Знак минус важен; на лекции студент сразу же его исправляет.

Есть и еще одно важное различие в области значений координаты. Для сферы новая координата пробегает от одного конца до другого и останавливается на \(\pi\). Здесь
\begin{equation}
0\leq \alpha<\infty.
\end{equation}
Обычный синус возрастает, а затем вновь обращается в нуль. Гиперболический синус вблизи \(\alpha=0\) начинается так же, но затем растет экспоненциально.

Это означает, что окружающие нас сферы меньшей размерности становятся всё больше и больше по мере продвижения наружу. В этой геометрии небо на всё больших радиусах не сначала увеличивается, а потом уменьшается. Оно продолжает расти.

В трех пространственных измерениях соответствующая форма FRW имеет вид
\begin{equation}
d\tau^2=dt^2-a(t)^2\,d\mathcal H_3^2,
\end{equation}
а метка знака такова:
\begin{equation}
k=-1.
\end{equation}

\section{Дисковое представление и его смысл}

Отрицательная кривизна становится интуитивно понятной в дисковом представлении, прославленном Эшером, особенно \emph{Limit Circle Number Four}. Аналогия со стереографической проекцией полезна, но поведение здесь противоположное.

Для сферы проекция выглядит бесконечной, тогда как реальная геометрия конечна. В случае отрицательной кривизны дисковая картина выглядит конечной, тогда как реальная геометрия бесконечна. Равные ``континенты'' или ``галактики'' кажутся всё меньше и меньше по мере приближения к границе, но геометрически все они одного и того же размера. Если измерять расстояние, подсчитывая, сколько одинаковых частей пересекается, то граница находится бесконечно далеко.

Таким образом, в пространстве отрицательной кривизны знакомый объект выглядит слишком маленьким по сравнению с плоским пространством. Он виден под слишком малым углом.

Геометрия по-прежнему однородна. Хотя на этой картине внешние части кажутся крошечными, любой наблюдатель, оказавшийся в другом месте, увидел бы ту же самую внутреннюю геометрию.

\begin{classroomqa}
\lecturetimestamp{01:06:42}
\audiencequestion{Наблюдатель в любой из этих окружностей увидел бы похожую картину?}
\lecturerresponse{Да, в точности ту же самую. Пространство однородно.}
\end{classroomqa}

\begin{classroomqa}
\lecturetimestamp{01:07:03}
\audiencequestion{В сферическом случае далекие вещи кажутся слишком большими, а здесь — слишком маленькими?}
\lecturerresponse{Да. При положительной кривизне знакомый объект виден под слишком большим углом по сравнению с плоским пространством. При отрицательной кривизне он виден под слишком малым углом.}
\end{classroomqa}

\begin{classroomqa}
\lecturetimestamp{01:08:09}
\audiencequestion{Связано ли значение стереографоподобной проекции с тем, что астроном может восстановить ее по наблюдениям?}
\lecturerresponse{Да. Астроном, по существу, — геодезист, измеряющий треугольники на небе. Для данного геометрического вопроса зависимость от времени можно оставить в стороне и сосредоточиться на углах и расстояниях.}
\end{classroomqa}

Есть и полезная локальная картина. Небольшой участок отрицательной кривизны можно изобразить как седловую поверхность в обычном трехмерном пространстве. Но таким способом можно изобразить только малые участки; всю гиперболическую плоскость нельзя глобально вложить как обычную поверхность в трехмерное евклидово пространство.

\section{Плоский случай}

Между случаями положительной и отрицательной кривизны находится плоское пространство. О нем можно думать как о предельном случае, но достаточно просто иметь в виду обычное плоское евклидово пространство.

Его форма FRW имеет вид
\begin{equation}
d\tau^2=dt^2-a(t)^2\left(dx^2+dy^2+dz^2\right),
\end{equation}
при
\begin{equation}
k=0.
\end{equation}

Вот три однородные геометрии FRW:
\begin{align}
k=+1 &: \text{positive spatial curvature},\\
k=0 &: \text{flat space},\\
k=-1 &: \text{negative spatial curvature}.
\end{align}

\begin{classroomqa}
\lecturetimestamp{01:15:07}
\audiencequestion{Если спроецировать геометрию на плоскость и проследить сигнал, будет ли проецированная скорость света меняться от места к месту?}
\lecturerresponse{Да, на проекции она менялась бы. На сферической проекции казалось бы, что свет движется всё быстрее и быстрее по мере удаления; на диске отрицательной кривизны казалось бы, что у края он движется всё медленнее и медленнее. Это изменение относится к картинке, а не к лежащей в основе физике.}
\end{classroomqa}

\begin{classroomqa}
\lecturetimestamp{01:17:05}
\audiencequestion{Действительно ли такой способ рассмотрения далеких объектов использовался для определения плоскостности?}
\lecturerresponse{Да, но на деле используется именно угловая информация. Плоскостность не определяют, пытаясь измерить, с какой скоростью свет, как кажется, движется на проекции.}
\end{classroomqa}

\section{Возвращение к космологическим уравнениям}

После перерыва появляется связь между уже изученными космологическими уравнениями и уравнениями Эйнштейна. В этот момент на доске всё еще остается фрагмент плоской пространственной метрики.

\begin{figure}[htbp]
\centering
\includegraphics[width=0.72\textwidth]{lecture_05_figure_02.png}
\caption{Фрагмент плоской метрики FRW.}
\lectureframe{01:18:18}
\end{figure}

Сначала идет очень краткий ньютоновский обзор. Поместим себя в центр. Рассмотрим галактику, удаляющуюся согласно закону Хаббла. По теореме Ньютона для сферически симметричного распределения вклад в силу, действующую на нее, дает только масса внутри сферы, радиус которой равен расстоянию до галактики. Тогда есть три возможности: галактика движется со скоростью выше скорости убегания, ниже скорости убегания или точно со скоростью убегания.

Далее напоминается, что для обычного нерелятивистского вещества плотность ведет себя как масса, деленная на \(a^3\), для излучения — как \(1/a^4\), а для энергии вакуума она не меняется при изменении сферы.

В более раннем ньютоновском рассмотрении космологическое уравнение содержало член с плотностью и, если не находиться точно в случае скорости убегания, дополнительный постоянный член, пропорциональный \(1/a^2\). В предыдущем обсуждении этот член по большей части полагали равным нулю. Эмпирически он мал, но концептуально нет причин его опускать.

\section{Уравнения Эйнштейна и компонента \(00\)}

Следующий шаг — уравнения поля Эйнштейна. В левой части стоит геометрия, построенная из кривизны пространства-времени. В правой части стоит тензор энергии-импульса, источник гравитационного поля. В нормировке, используемой в лекции, пишут
\begin{equation}
R^{\mu\nu}-\frac{1}{2}g^{\mu\nu}R=4\pi G\,T^{\mu\nu}.
\end{equation}
\editorialnote{В видимой на доске строке коэффициент \(\tfrac12\) неразборчив, но около 01:33:36 устное исправление его восстанавливает. Здесь это уравнение сохраняет это исправление при сохранении лекционной нормировки \(4\pi G\).}

Здесь нужны не все десять независимых компонент. К этому моменту уже предположено так много --- однородность, форма FRW и единственный зависящий от времени масштабный фактор --- что остается только одна произвольная функция, а именно \(a(t)\). Естественно выбрать компоненту время-время, потому что \(T^{00}\) — это плотность энергии \(\rho\).

\begin{figure}[htbp]
\centering
\includegraphics[width=0.62\textwidth]{lecture_05_figure_05.png}
\caption{Уравнение поля Эйнштейна и его компонента \(00\).}
\lectureframe{01:30:54}
\end{figure}

Итак, в центре внимания стоит
\begin{equation}
R^{00}-\frac{1}{2}g^{00}R=4\pi G\,\rho.
\end{equation}

Правая часть — это плотность энергии. Левая часть — геометрия.

\begin{classroomqa}
\lecturetimestamp{01:39:09}
\audiencequestion{А как насчет других компонент тензора энергии-импульса?}
\lecturerresponse{Вообще говоря, они не обращаются в нуль. Пространственно-пространственные компоненты могут представлять давление. В случае доминирования вещества давление равно нулю, в случае доминирования излучения оно составляет одну треть плотности энергии, а для энергии вакуума пишут \(P=w\rho\) при \(w=-1\).}
\end{classroomqa}

\begin{classroomqa}
\lecturetimestamp{01:40:24}
\audiencequestion{Значит, если бы было только вещество, остальные компоненты \(T^{\mu\nu}\) обращались бы в нуль?}
\lecturerresponse{Для вещества без давления — да, компоненты давления обращаются в нуль.}
\end{classroomqa}

\section{Что дает вычисление по Эйнштейну}

Полный расчет кривизны на доске не проводится. Важно то, что получается, когда метрику FRW подставляют в тензор Эйнштейна.

Один тип члена возникает просто из того факта, что масштабный фактор зависит от времени. Даже если бы пространство было плоским, зависящий от времени масштабный фактор всё равно давал бы кривизну пространства-времени. Этот вклад пропорционален
\begin{equation}
\left(\frac{\dot a}{a}\right)^2.
\end{equation}

Второй тип члена возникает из самой пространственной кривизны. Даже если бы масштабный фактор не менялся со временем, искривлённое пространство всё равно вносило бы вклад в кривизну пространства-времени. Для простой геометрии постоянной кривизны кривизна ведёт себя как величина, обратная квадрату радиуса кривизны, так что она даёт член, пропорциональный
\begin{equation}
\frac{k}{a^2},
\end{equation}
со знаком плюс для сферы, со знаком минус для гиперболического случая и равный нулю для плоского пространства.

\begin{classroomqa}
\lecturetimestamp{01:36:53}
\audiencequestion{Не должен ли здесь также присутствовать такой множитель, как радиус кривизны?}
\lecturerresponse{Это уже и есть множитель \(1/a^2\). Для этих простых геометрий постоянной кривизны кривизна — это просто обратный квадрат радиуса кривизны.}
\end{classroomqa}

Учёт коэффициентов ведётся по ходу вывода и частично проверяется сравнением с полученным ранее ньютоновским ответом. В конечном итоге получаем\footnote{Эта выписанная форма уравнения Фридмана объединяет устное обсуждение коэффициентов с исправлениями по ходу изложения и сравнение с полученным ранее ньютоновским результатом на отрезке 01:32:35--01:36:53.}
\begin{equation}
\left(\frac{\dot a}{a}\right)^2+\frac{k}{a^2}=\frac{8\pi G}{3}\rho,
\end{equation}
или, эквивалентно,
\begin{equation}
\left(\frac{\dot a}{a}\right)^2=\frac{8\pi G}{3}\rho-\frac{k}{a^2}.
\end{equation}

Это именно та форма, которая связывает три геометрии FRW с рассмотренной ранее космологической динамикой.

\begin{classroomqa}
\lecturetimestamp{01:37:52}
\audiencequestion{Может ли \(k\) принимать любое значение, если переопределить \(a\)?}
\lecturerresponse{В плоском случае для \(a\) нет единственного естественного масштаба, так что его можно перемасштабировать, но \(k\) остаётся нулём. В сферическом случае и в случае отрицательной кривизны существует естественный масштаб кривизны. Если взять \(a(t)\) равным этому естественному радиусу кривизны, то \(k\) будет просто \(+1\) или \(-1\).}
\end{classroomqa}

\begin{classroomqa}
\lecturetimestamp{01:41:00}
\audiencequestion{Равно ли \(k\) в том последнем уравнении \(+1\), \(-1\) или \(0\)?}
\lecturerresponse{Да. Для сферы это \(+1\), для гиперболического случая \(-1\), а для плоского пространства \(0\).}
\end{classroomqa}

Ньютоновское рассмотрение не ввело нас в заблуждение; оно дало то же основное уравнение. То, что добавляет общая теория относительности, — это интерпретация. С ньютоновской точки зрения дополнительная константа различает случаи выше скорости убегания, ниже неё и пограничный случай. С точки зрения Эйнштейна тот же знак говорит нам о геометрии самого пространства.

\begin{classroomqa}
\lecturetimestamp{01:42:46}
\audiencequestion{Как это работало до включения космологической постоянной?}
\lecturerresponse{До введения космологической постоянной между режимом относительно скорости убегания и геометрией пространства существовала тесная связь. Если \(k=+1\), пространство имеет положительную кривизну, а движение соответствует случаю ниже скорости убегания, так что Вселенная может расширяться, а затем вновь сжиматься. Если \(k=-1\), пространство имеет отрицательную кривизну, а движение соответствует случаю выше скорости убегания, так что Вселенная продолжает расширяться. Если \(k=0\), имеет место предельный, пограничный случай.}
\end{classroomqa}

Следовательно, один и тот же знак \(k\) можно читать двумя способами: как динамическую классификацию движения и как знак кривизны пространства.
